\lecture{14}{2026-10-09}{Algebraic Power Series}

Given \(P(z, y) = P_d(z) y^d + P_{d-1}(z) y^{d-1} + \cdots + P_0(z) \in \mathbb{Q}[z][y]\) with \(P_d(z) \neq 0\).
Initial terms \(F(z) = f_0 + f_1 z + f_2 z^2 + \cdots\) of a solution \(y = F(z)\) of \(P(z, y) = 0\) are enough to specify a unique solution \(F(z)\) in \(\mathbb{Q}[[z]]\).

The goal is to find an asymptotic expansion of \((f_n)\).
How do we do that directly from \(P(z, y)\)?

Recall \(\mathbb{Q}[[z]]\) is the field of formal power series,
\[
    \mathbb{Q}((z)) = \left\{ \sum_{n \geq N} a_n z^n : N \in \mathbb{Z}, a_n \in \mathbb{Q} \right\}
\]
is the field of formal Laurent series, and
\[
    \mathbb{Q}^{\mathrm{fra}}((z)) = \left\{ \sum_{n \geq N} a_n z^{n/R} : N \in \mathbb{Z}, R \in \mathbb{Z}_{>0}, a_n \in \mathbb{Q} \right\}
\]
is the field of formal fractional power series.

\begin{theorem}[Puiseux's Theorem]
    \(\mathbb{Q}^{\mathrm{fra}}((z))\) is the algebraic closure of \(\mathbb{Q}((z))\).
\end{theorem}

\begin{example}[Non-example]
    The series \(\sum_{n \geq 1} z^{1/n}\) is not in \(\mathbb{Q}^{\mathrm{fra}}((z))\).
\end{example}

We say \(F(z)\) is \emph{algebraic} if there exists a non-zero polynomial \(P(z, y) \in \mathbb{Q}[z][y]\) such that \(P(z, F(z)) = 0\).
A \emph{minimal polynomial} of \(F(z)\) is a polynomial \(P(z, y) \in \mathbb{Q}[z][y]\) such that \(P(z, F(z)) = 0\) and \(\deg_y P\) is minimal among all such polynomials.
The other Puiseux series solutions of \(P(z, y) = 0\) are called \emph{conjugates} of \(F(z)\).

There is an algorithm (Newton-Puiseux algorithm) to find the Puiseux series solutions of \(P(z, y) = 0\).

\begin{example}
    Consider
    \[
        P(z, y) = zy^2 - y + 1.
    \]
    Puiseux's Theorem implies that there exist solutions of the form
    \[
        F(z) = a_r z^r + \cdots
    \]
    for some \(r \in \mathbb{Q}\) and \(a_r \neq 0\), where \(\cdots\) denotes higher order terms in \(z\).
    Plugging this into \(P(z, F(z)) = 0\) gives
    \[
        z(a_r z^r + \cdots)^2 - (a_r z^r + \cdots) + 1 = 0.
    \]
    Let's compare the lowest order terms above.
    The lowest order terms of each summand are \(a_r^2 z^{2r + 1}\), \(-a_r z^r\), and \(1\), respectively.
    Since they must sum to zero, at least two of them must have the same order.
    \begin{description}
        \item[\boldmath If \(2r + 1 = r \leq 0\),] then \(r = -1\) and \(a_r^2 - a_r = 0\), so \(a_r = 1\).
        \item[\boldmath If \(2r + 1 = 0 \leq r\),] then \(r = -\frac{1}{2}\), however \(0 \not\leq -\frac{1}{2}\), so this case is impossible.
        \item[\boldmath If \(r = 0 \leq 2r + 1\),] then \(r = 0\) and \(-a_r + 1 = 0\), so \(a_r = 1\).
    \end{description}
    Two plausible roots
    \[
        f_1(z) = 1 + \cdots, \quad f_2(z) = z^{-1} + \cdots.
    \]
    To find more terms, make a substitution and do this again.
    For example, write \(f_1(z) = 1 + z g(z)\), where the lowest order term of \(g(z)\) is of order \(> -1\).
    Then \(g(z)\) is a solution of
    \[
        Q(z, y) = P(z, 1 + z y) = z^2 y^2 + z y.
    \]
    We can now apply the same method to find the lowest order term of \(g(z)\).
\end{example}

A student asked if this is how it's computed in practice.
This is particularly slow.
You can use it to figure out enough lowest terms to separate the roots, then you can use Newton's method to find the rest of the terms in a more efficient way.

The summary of the Newton's method is as follows.
Suppose \(F(0) = 0\), so that \(P(0, 0) = 0\), and suppose \(P_y(0, 0) \neq 0\).
Define \(F_0(z) = 0\) and
\[
    F_{n+1}(z) = F_n(z) - \frac{P(z, F_n(z))}{P_y(z, F_n(z))} \mod z^{2^{n+1}}
\]
Then, \(F_{n+1}(z)\) agrees with \(F_n(z)\) up to order \(2^{n+1}\).

Let \(P(z, y) \in \mathbb{Z}[z, y]\) (scale if needed to make it integral) and suppose \(f(z) = \sum_{n \geq N} c_n z^{n/R}\) satisfies \(P(z, f(z)) = 0\).

The \emph{singular part} of \(f\) is the initial partial sum
\[
    f_\ell(z) = \sum_{n = N}^{\ell} c_n z^{n/R}
\]
where \(\ell \geq N\) is the minimum integer such that \(f_\ell(z)\) is not equal to the initial partial sum of any other root of \(P(z, y) = 0\).

The \emph{separation order} of \(f\) with respect to \(P\) is \(\ell\).

\begin{proposition}
    Let \(P(z, y) \in \mathbb{Z}[z, y]\) be irreducible in \(y\) 
\end{proposition}